\section*{第 \ref{ch:b2:ligand-field} 章——配体场理论与颜色}

\begin{solution}{exo:b2:ligand-field:1}
$d^3$：$t_{2g}^3$，\omterm{def:b2:ligand-field:cfse}{晶体场稳定化能} $-1.2\Delta_o$。$d^8$：$t_{2g}^6e_g^2$，\omterm{def:b2:ligand-field:cfse}{晶体场稳定化能} $-2.4 + 1.2 = -1.2\Delta_o$。
\end{solution}

\begin{solution}{exo:b2:ligand-field:2}
\ce{[Fe(H2O)6]^2+}，$d^6$ \omterm{def:b2:ligand-field:spin}{高自旋} $t_{2g}^4e_g^2$：四个。\ce{[Fe(CN)6]^4-}，\omterm{def:b2:ligand-field:spin}{低自旋}
$t_{2g}^6$：无（\omterm{def:b2:diatomic-mos:magnetism}{抗磁性}）。
\end{solution}

\begin{solution}{exo:b2:ligand-field:3}
600 nm 是橙色光；该配合物呈蓝色（青蓝色）。
\end{solution}

\begin{solution}{exo:b2:ligand-field:4}
$\ce{Cl-} < \ce{H2O} < \ce{NH3} < \ce{CN-}$.
\end{solution}

\begin{solution}{exo:b2:ligand-field:5}
水：$\Delta_o < P$，\omterm{def:b2:ligand-field:spin}{高自旋}，$t_{2g}^3e_g^2$，五个未成对电子。氰根：$\Delta_o > P$，
\omterm{def:b2:ligand-field:spin}{低自旋}，$t_{2g}^5$，一个未成对电子。
\end{solution}

\begin{solution}{exo:b2:ligand-field:6}
在四面体型 \ce{[CoCl4]^2-} 中，\omterm{def:b2:ligand-field:crystal-field}{晶体场分裂}约为八面体情形的 $4/9$，而
氯离子是比水更弱的配体：d–d 跃迁吸收带向长波方向移动（由橙色
移至红色），该配合物呈蓝色。四面体配合物的吸收也更强，
因此颜色更深。
\end{solution}

\begin{solution}{exo:b2:ligand-field:7}
\ce{Ni^2+} 为 $d^8$。平面正方形：四个较低的轨道以成对方式容纳八个电子，
$\mathrm d_{x^2-y^2}$ 轨道保持空置：\omterm{def:b2:diatomic-mos:magnetism}{抗磁性}。四面体构型：$e^4t_2^4$，两个未成对电子
在 $t_2$ 组中。
\end{solution}

\begin{solution}{exo:b2:ligand-field:8}
$\tilde\nu = 1/(500 \times 10^{-7}\,\text{cm}) = \qty{20000}{cm^{-1}}$; $\Delta_o =
0.1196/(500 \times 10^{-9}) = \qty{2.39e5}{J/mol}$, \qty{239}{kJ/mol}.
\end{solution}

\begin{solution}{exo:b2:ligand-field:9}
\ce{Zn^2+} 为 $d^{10}$，\ce{Sc^3+} 为 $d^0$，因此不存在 d–d 跃迁
（\ce{Zn^2+} 无空的 d 能级，\ce{Sc^3+} 无 d 电子），可见光区无吸收。
\end{solution}

\begin{solution}{exo:b2:ligand-field:10}
Co(III) 为 $d^6$，加上六个氨的孤电子对，共 18 个电子。十二个电子填满六个\omterm{def:b2:diatomic-mos:bonding}{成键轨道}，
六个电子填入非键的 $t_{2g}$ 轨道，氨给出的场对钴足够强，
形成\omterm{def:b2:ligand-field:spin}{低自旋}；$e_g^*$ 轨道为空。所有电子均已成对，呈\omterm{def:b2:diatomic-mos:magnetism}{抗磁性}。
\end{solution}

\begin{solution}{exo:b2:ligand-field:11}
\ce{CO} 是良好的 $\sigma$ 给体，而且更重要的是 $\pi$ 受体：其空的 $\pi^*$ 轨道
通过\omterm{def:b2:ligand-field:pi-ligands}{反馈键合}降低 $t_{2g}$ 能级，使 $\Delta_o$ 增大。\ce{F-} 是 $\pi$ 给体，这使
$t_{2g}$ 能级升高，$\Delta_o$ 减小。点电荷模型只看到电荷，于是
排序就错了。
\end{solution}

\begin{solution}{exo:b2:ligand-field:12}
若无晶体场，水合焓本应沿一条平滑曲线变化（离子沿该系列
半径逐渐减小）。水给出\omterm{def:b2:ligand-field:spin}{高自旋}的水合配合物，对 $d^0$、$d^5$、$d^{10}$ 而言其\omterm{def:b2:ligand-field:cfse}{晶体场稳定化能}为零，
而对 $d^3$ 和 $d^8$ 则最大，额外的稳定化使曲线出现两个凸起，峰值
靠近 \ce{V^2+} 和 \ce{Ni^2+}，高于经过 \ce{Ca^2+}、\ce{Mn^2+}、\ce{Zn^2+} 的直线。
\end{solution}

\begin{solution}{pb:b2:ligand-field:1}
\textbf{1.} $d^3$.
\textbf{2.} $t_{2g}^3$，三个轨道中各有一个电子，自旋平行。
\textbf{3.} 否：三个电子无需成对即可填入 $t_{2g}$。
\textbf{4.} $-1.2\Delta_o$.
\textbf{5.} 三个。
\textbf{6.} \omterm{def:b2:ligand-field:cfse}{晶体场稳定化能}较大，且指向配体的 $e_g^*$ 轨道
为空：移去或加入配体都会损失相当一部分该稳定化能。
\textbf{7.} 六个氧离子位于一个略有畸变的八面体的顶点上。
\textbf{8.} $1/(556 \times 10^{-7}) = \qty{17990}{cm^{-1}}$，约 \qty{1.80e4}{cm^{-1}}。
\textbf{9.} $0.1196/(556 \times 10^{-9}) = \qty{2.15e5}{J/mol}$: \qty{215}{kJ/mol}.
\textbf{10.} 黄绿色（556 nm）和紫色（405 nm）。
\textbf{11.} 红色，并带一点蓝色：即红宝石的深红色。
\textbf{12.} $d^3$ 离子有多个激发组态，其中一个电子位于 $e_g$；
电子间的排斥使它们能量不同，因此出现多个吸收带（在
本科三年级分册中讨论）。
\textbf{13.} $0.1196/(620 \times 10^{-9}) = \qty{1.93e5}{J/mol}$, \qty{193}{kJ/mol}
(\qty{16130}{cm^{-1}}).
\textbf{14.} 红宝石：其\omterm{def:b2:ligand-field:crystal-field}{晶体场分裂}更大。
\textbf{15.} 变为橙红色：红光现在被吸收，绿色（带一些蓝色）被透过。
\textbf{16.} 在绿柱石中，铬所在的位置略大，氧离子距离更远，
周围环境也不同，场较弱。
\textbf{17.} 否：$d^3$ 离子在任何八面体场中都保留三个未成对电子。
\textbf{18.} $d^9$.
\textbf{19.} 水合配合物通过 d–d 跃迁吸收红色和橙色光：
透射光呈蓝色。
\textbf{20.} 溶液变为白色：铜周围没有水配体（硫酸根占据了它们的
位置），场变弱，吸收带移入红外区。
\textbf{21.} 水重新与铜结合，形成蓝色的水合配合物。
\textbf{22.} 吸收发生在更短的波长：氨在\omterm{def:b2:ligand-field:spectrochemical}{光谱化学序列}中排在水之后，给出更大的
分裂（即氨合配合物的深蓝紫色）。
\textbf{23.} 否：$t_{2g}^6e_g^3$ 是唯一的组态。
\textbf{24.} 是：一个未成对电子。
\textbf{25.} $\boldsymbol{\Delta_o \approx \qty{1.8e4}{cm^{-1}}}$，约
$\boldsymbol{\qty{215}{kJ/mol}}$.
\end{solution}
